\section{Fine moduli spaces of modules}
\label{subsec:finemod}

\subsection{Cornering and fine moduli spaces}

Let $\Rep(\Gamma)$
denote the representation ring of $\Gamma$. Introduce a formal symbol $\rho_{\infty}$ so that $\{\rho_i \mid i \in Q_0\}$ provides a $\mathbb{Z}$-basis for $\mathbb{Z}^{Q_0}\cong \mathbb{Z}\oplus \Rep(\Gamma)$ considered as $\mathbb{Z}$-modules.

Recall that $\PiQ=\Pi/(b^{\ast})$ is the quotient of the
preprojective algebra of the framed McKay quiver of $\Gamma$ by the two sided ideal $(b^{\ast})$, where $b^\ast$ is the arrow from $0$ to $\infty$.
Let $I \subseteq \{0, \dots , r\}$ be a non-empty subset; we work under the assumption that $I\ne \varnothing$ for the rest of the paper.
Define the idempotent $\overline{e}_I \coloneqq e_{\infty} + \sum_{i \in I} e_i \in \PiQ$.
The subalgebra
\begin{equation}\label{eqn:AI}
\PiQ_I \coloneqq \overline{e}_I \PiQ \overline{e}_I
\end{equation}
of $A$ comprises linear combinations of
classes of paths in $Q$ whose tails and heads lie in the set $\{\infty\} \cup I$; this is also an instance of \emph{cornering}~\cite[Remark~3.1]{craw2018multigraded}.

Let
$n_I \coloneqq {(n_i)_{i\in I}}= \sum_{i \in I} n_i \rho_i \in \mathbb{N}^{I}$.
Then $(1,n_I)\coloneqq\rho_{\infty} +n_I \in \mathbb{N} \oplus \mathbb{N}^{I}$
is a dimension vector for $\PiQ_I$-modules, and we consider the stability condition $\eta_{I}\colon \mathbb{Z} \oplus \mathbb{Z}^{I} \to \mathbb{Q}$
given by
\begin{equation}
 \label{eqn:etaI}
\eta_I (\rho_i) =
\begin{cases}
\displaystyle -\sum_{j \in I} n_j & \textrm{for } i= \infty, \\
1 & \textrm{if }i \in I.
\end{cases}
\end{equation}

It follows directly from the definition that an $\PiQ_I$-module $N$ of dimension vector $(1,n_I)$ is $\eta_I$-stable if and only if
there exists a surjective $\PiQ_I$-module homomorphism $\PiQ_I e_{\infty} \to N$.
The quiver moduli space construction of King~\cite[Proposition~5.3]{king1994moduli} for finite dimensional algebras can be adapted to any algebra presented as the quotient of a finite, connected quiver by an ideal of relations. We established in \cite[Proposition~3.3]{craw2019punctual} that the algebra $\PiQ_I$ has this property.
The vector~$(1,n_I)$ is indivisible and $\eta_I$ is a generic stability condition for $\PiQ_I$-modules of dimension vector~$(1,n_I)$, so there is a~fine moduli space $\mathcal{M}_{\PiQ_I}(1,n_I)$
of $\eta_I$-stable $\PiQ_I$-modules of dimension
vector~$(1,n_I)$. Let $T_I \coloneqq\bigoplus_{i \in \{\infty \}\cup I } T_i$ denote the tautological vector bundle on~$\mathcal{M}_{\PiQ_I}(1,n_I)$, where~$T_{\infty}$ is the trivial bundle and $T_i$ has rank~$n_i$ for $i \in I$.
After multiplying $\eta_I$ by a positive integer $m\geq 1$ if necessary, we may assume that the polarising ample line bundle $\mathcal{L}_I \coloneqq \bigotimes_{i \in I} \mathrm{det}(T_i)^{m}$ on $\mathcal{M}_{\PiQ_I}(1,n_I)$ given by the GIT construction is very ample.

